%======================================================================
% Section 9: Reflection, Transmission, and Interface Phenomena
%======================================================================
\section{Reflection, Transmission, and Interface Phenomena}
\label{sec:interfaces}

This section presents wave reflection, transmission, and energy partitioning at interfaces between different continuum media. The scattering matrix formulation, interface conditions, and energy ratios are developed for elastic, thermoelastic, and piezothermoelastic media.

\subsection{Interface geometry}
\label{sec:9.1}

Consider two half-spaces in perfect contact at the interface plane $x_3 = 0$. Medium 1 occupies $x_3 < 0$ and medium 2 occupies $x_3 > 0$. A plane harmonic wave is incident from medium 1 onto the interface, generating reflected waves in medium 1 and transmitted waves in medium 2.

The interface normal is $\bm{n} = (0, 0, 1)$. All field variables vary harmonically in time with convention $e^{-i\omega t}$ and propagate along the interface with wave number $k$ in the $x_1$-direction:

\begin{equation}
\text{field}(x_1, x_3, t) = \hat{\text{field}}(x_3) e^{i(kx_1 - \omega t)}.
\label{eq:9.1}
\end{equation}

\subsection{Wave decomposition: incident, reflected, transmitted modes}
\label{sec:9.2}

In each medium, the general solution is a superposition of partial waves. For an isotropic elastic medium, there are two wave types: longitudinal (P) and shear (SV). For thermoelastic media, there is an additional thermal (T) wave. For piezothermoelastic media, there are additional coupled modes (qP, qS, qT, electromagnetic).

The displacement field in medium $m$ ($m = 1, 2$) is decomposed as:

\begin{equation}
u_i^{(m)} = \sum_n A_n^{(m)} \hat{u}_i^{(m,n)} e^{i(q_n^{(m)}x_3 - \omega t)},
\label{eq:9.2}
\end{equation}

where $A_n^{(m)}$ are complex amplitude coefficients, $q_n^{(m)}$ are the vertical wave numbers, and $\hat{u}_i^{(m,n)}$ are the polarization vectors.

The incident wave comes from medium 1 with known amplitude. Reflected waves propagate back into medium 1, and transmitted waves propagate into medium 2.

\subsection{Interface conditions}
\label{sec:9.3}

At a perfectly bonded interface between two media, the following conditions must hold:

\subsubsection{Mechanical conditions}

\begin{equation}
u_i^{(1)} = u_i^{(2)}, \qquad \sigma_{j3}^{(1)} = \sigma_{j3}^{(2)} \quad \text{at } x_3 = 0.
\label{eq:9.3}
\end{equation}

The first condition enforces displacement continuity (no slip or separation); the second enforces traction equilibrium (Newton's third law at the interface).

\subsubsection{Thermal conditions}

For thermoelastic media:

\begin{equation}
T^{(1)} = T^{(2)}, \qquad q_3^{(1)} = q_3^{(2)} \quad \text{at } x_3 = 0.
\label{eq:9.4}
\end{equation}

Temperature continuity and heat flux continuity ensure perfect thermal contact.

\subsubsection{Electrical conditions (piezoelectric media)}

For piezothermoelastic media:

\begin{equation}
\phi^{(1)} = \phi^{(2)}, \qquad D_3^{(1)} = D_3^{(2)} \quad \text{at } x_3 = 0.
\label{eq:9.5}
\end{equation}

Electric potential continuity and electric displacement continuity ensure perfect electrical contact. For an insulated interface, $D_3 = 0$ replaces the second condition.

\subsubsection{Imperfect bonding}

For imperfectly bonded interfaces, spring-type conditions relate displacement jumps to interfacial tractions:

\begin{equation}
u_i^{(2)} - u_i^{(1)} = K_i^u t_i, \qquad \sigma_{j3}^{(1)}n_j = \sigma_{j3}^{(2)}n_j = t_i,
\label{eq:9.6}
\end{equation}

where $K_i^u$ is the interfacial stiffness tensor and $t_i$ is the interfacial traction. In the limit $K_i^u \to \infty$, perfect bonding is recovered; in the limit $K_i^u \to 0$, a traction-free (slip) interface is recovered.

\subsection{Scattering matrix formulation}
\label{sec:9.4}

The interface conditions lead to a linear system relating the unknown amplitude coefficients to the known incident wave amplitude. For a thermo-piezoelastic interface problem with $N$ unknown amplitudes and $N$ interface conditions, the system can be written in matrix form:

\begin{equation}
\Delta \bm{X} = \bm{Q},
\label{eq:9.7}
\end{equation}

where:
\begin{itemize}
\item $\Delta$ is the $N\times N$ system matrix whose entries depend on material properties, wave numbers, frequency, and incidence angle,
\item $\bm{X} = (A_1, A_2, \ldots, A_N)^T$ is the vector of unknown amplitude coefficients (reflected and transmitted),
\item $\bm{Q} = (Q_1, Q_2, \ldots, Q_N)^T$ is the known vector determined by the incident wave amplitude.
\end{itemize}

The entries of $\Delta$ are determined by substituting the partial wave expansions into the interface conditions \eqref{eq:9.3}--\eqref{eq:9.5}. Each row corresponds to one interface condition; each column corresponds to one unknown amplitude.

The solution is obtained by matrix inversion:

\begin{equation}
\bm{X} = \Delta^{-1}\bm{Q}.
\label{eq:9.8}
\end{equation}

The solution exists and is unique when $\det(\Delta) \neq 0$.

\subsection{Secular equation and mode classification}
\label{sec:9.5}

The secular equation (characteristic polynomial) determines the allowed wave numbers in each medium:

\begin{equation}
G_1 q^8 + G_2 q^6 + G_3 q^4 + G_4 q^2 + G_5 = 0,
\label{eq:9.9}
\end{equation}

which is an eighth-order polynomial in $q$ for coupled thermo-piezoelastic media. The roots are classified as:

\begin{itemize}
\item \textbf{Propagating modes:} real $q$ --- waves that propagate without spatial attenuation,
\item \textbf{Evanescent modes:} purely imaginary $q$ --- waves that decay exponentially with distance,
\item \textbf{Attenuating modes:} complex $q$ --- waves that both propagate and attenuate.
\end{itemize}

Only outgoing modes (satisfying the radiation condition) are retained in the physical solution.

\subsection{Number of boundary conditions vs. unknown amplitudes}
\label{sec:9.6}

For a well-posed interface problem, the number of independent boundary conditions must equal the number of unknown amplitude coefficients. Table \ref{tab:T6_interface} summarizes interface condition configurations for different interface types.

\begin{figure}[htbp]
\centering
\includegraphics[width=0.9\textwidth]{fig03-interface_geometry.pdf}
\caption{Interface reflection and transmission geometry: incident P-wave from Medium 1 generates reflected P and S waves in Medium 1 and transmitted qP and qS waves in Medium 2. Interface conditions enforce displacement, traction, temperature, heat flux, electric potential, and electric displacement continuity.}
\label{fig:interface_geometry}
\end{figure}

%======================================================================
% Table T6: Interface configurations
%======================================================================
\begin{table}[htbp]
\centering
\small
\caption{Interface condition configurations for wave reflection and transmission problems.}
\label{tab:T6_interface}
\begin{tabularx}{\textwidth}{@{}lXXXXX@{}}
\toprule
\textbf{Interface type} & \textbf{Mechanical conditions} & \textbf{Thermal conditions} & \textbf{Electrical conditions} & \textbf{Applications} \\
\midrule
Perfectly bonded (welded) & $u_i^{(1)} = u_i^{(2)}$; $\sigma_{j3}^{(1)} = \sigma_{j3}^{(2)}$ & $T^{(1)} = T^{(2)}$; $q_3^{(1)} = q_3^{(2)}$ & N/A (non-piezoelectric) & Welded joints, composite laminates, geological interfaces \\
Perfectly bonded + piezoelectric & Same mechanical & Same thermal & $\phi^{(1)} = \phi^{(2)}$; $D_3^{(1)} = D_3^{(2)}$ & Piezoelectric sensors, transducers, smart structures \\
Insulated (thermally) & Same mechanical & $q_3 = 0$ & N/A & Thermal barriers, insulated boundaries \\
Insulated (electrically) & Same mechanical & Same thermal & $D_3 = 0$ & Electrically insulated piezoelectric interfaces \\
Imperfect bonding (spring-type) & $u_i^{(2)} - u_i^{(1)} = K_i^u t_i$; $\sigma_{j3}^{(1)} = \sigma_{j3}^{(2)}$ & Same thermal & Same electrical & Adhesive layers, micro-cracks, weakly bonded interfaces \\
Thermal contact resistance & Same mechanical & $q_3^{(1)} = q_3^{(2)}$; $T^{(1)} - T^{(2)} = R q_3$ & N/A & Thermal interface materials, electronic cooling \\
\bottomrule
\end{tabularx}
\end{table}


\subsection{Energy ratios at interfaces}
\label{sec:9.7}

The energy reflection coefficient $R$, transmission coefficient $T$, and absorption coefficient $A$ are defined as:

\begin{equation}
R = \frac{|\avg{J_{\mathrm{ref},n}}_t|}{|\avg{J_{\mathrm{inc},n}}_t|},
\qquad
T = \frac{|\avg{J_{\mathrm{trans},n}}_t|}{|\avg{J_{\mathrm{inc},n}}_t|},
\qquad
A = \frac{\mathcal{D}}{|\avg{J_{\mathrm{inc},n}}_t|},
\label{eq:9.10}
\end{equation}

where the averages are time averages over one period and $n$ is the interface normal.

For a conservative system with no dissipation ($\mathcal{D} = 0$) and a passive interface:

\begin{equation}
\sum R + \sum T = 1.
\label{eq:9.11}
\end{equation}

With dissipation:

\begin{equation}
\sum R + \sum T + A = 1.
\label{eq:9.12}
\end{equation}

The identity \eqref{eq:9.12} holds when the energy flux is correctly defined for the specific theory, the interface is passive, all modes are properly accounted for, and complex modes are consistently normalized.

\subsection{Notation summary for Section 9}

\begin{table}[htbp]
\centering
\small
\caption{Notation introduced in Section 9.}
\label{tab:notation-sec9}
\begin{tabularx}{\textwidth}{@{}lXlr@{}}
\toprule
\textbf{Symbol} & \textbf{Meaning} & \textbf{Units} & \textbf{Equation} \\
\midrule
$A_n^{(m)}$ & Complex amplitude coefficient & varies & \eqref{eq:9.2} \\
$q_n^{(m)}$ & Vertical wave number & 1/m & \eqref{eq:9.2} \\
$\Delta$ & System matrix & varies & \eqref{eq:9.7} \\
$\bm{X}$ & Unknown amplitude vector & varies & \eqref{eq:9.7} \\
$\bm{Q}$ & Known incident vector & varies & \eqref{eq:9.7} \\
$R, T, A$ & Reflection, transmission, absorption coefficients & dimensionless & \eqref{eq:9.10} \\
$K_i^u$ & Interfacial stiffness & varies & \eqref{eq:9.6} \\
\bottomrule
\end{tabularx}
\end{table}

\subsection{Open questions for Section 9}

\begin{enumerate}[label=\textbf{OQ9.\arabic*}]
\item What are the interface conditions for nonlocal media where stress depends on strain in a neighborhood? \\ \label{oq:9.1}
\item How are energy ratios computed for memory-dependent media with hereditary flux? \\ \label{oq:9.2}
\item What is the correct interface energy balance for piezothermoelastic media with electrical coupling? \\ \label{oq:9.3}
\end{enumerate}

% LocalWords: 
